\section{Optimasi dan Metaheuristik untuk Konfigurasi ABAC Edge}

Konfigurasi ABAC berbasis \textit{edge} dalam penelitian ini dipandang sebagai masalah optimasi multi-parameter karena melibatkan pemilihan atribut yang boleh di-\textit{cache}, penentuan TTL per atribut, dan ukuran \textit{cache}. \textcite{weerasinghe_adaptive_2023} menunjukkan bahwa \textit{adaptive context caching} pada aplikasi IoT perlu mengoptimalkan performa dan efisiensi biaya ketika sistem merespons \textit{context query} secara \textit{real-time}. Selain itu, penelitian tersebut juga menekankan bahwa \textit{adaptive caching} mencakup proses \textit{selection, refreshing, scaling}, dan \textit{eviction} yang saling memengaruhi dalam manajemen siklus hidup \textit{cache}. \textcite{medvedev_refresh_2023} menjelaskan bahwa \textit{freshness, accuracy, completeness,} dan latensi merupakan parameter QoS, sedangkan biaya operasional perlu dipertimbangkan secara bersamaan dalam pengelolaan \textit{context caching} berbasis IoT \textit{middleware}. Dengan demikian, konfigurasi \textit{cache} dan TTL tidak cukup ditentukan secara manual karena setiap parameter dapat memengaruhi latensi, \textit{PIP calls, cache hit,} dan risiko penggunaan atribut yang usang.

\textit{Random search} digunakan sebagai \textit{baseline} optimasi karena pendekatan ini memilih konfigurasi secara acak dari ruang solusi yang telah ditentukan. \textcite{bergstra_random_2012} menunjukkan bahwa \textit{random search} dapat menjadi strategi yang lebih efisien dibandingkan \textit{grid search} ketika hanya sebagian parameter yang memiliki pengaruh besar terhadap performa sistem. Selain itu, \textcite{feng_dynamic_2025} menggunakan \textit{random search} sebagai \textit{baseline} pembanding pada optimasi \textit{multi-objective} berbasis GWO, bersama \textit{Genetic Algorithm} (GA). Oleh karena itu, \textit{random search} relevan digunakan sebagai \textit{baseline} awal untuk mengevaluasi apakah algoritma metaheuristik mampu menemukan konfigurasi TTL, ukuran \textit{cache}, dan \textit{cacheable attribute} yang lebih optimal dibandingkan pemilihan konfigurasi secara acak.

GWO digunakan sebagai metode utama karena mampu melakukan pencarian berbasis \textit{fitness} pada ruang solusi yang kompleks tanpa membutuhkan model analitik tertutup. \textcite{mirjalili_grey_2014} menjelaskan bahwa GWO meniru hirarki sosial serigala abu-abu yang terdiri dari \textit{alpha, beta, delta,} dan \textit{omega}, di mana kandidat solusi terbaik memandu proses pembaruan posisi kandidat lain. Dalam GWO, proses optimasi dilakukan secara \textit{derivation-free} dan berbasis pencarian stokastik sehingga cocok digunakan pada permasalahan optimasi dengan ruang pencarian yang kompleks dan sulit dimodelkan secara matematis. Dalam penelitian ini, mekanisme tersebut relevan karena setiap kandidat solusi dapat merepresentasikan kombinasi konfigurasi \textit{cacheable attribute}, TTL per atribut, dan ukuran \textit{cache}, sedangkan \textit{fitness} dihitung berdasarkan gabungan metrik performa dan keamanan.

Relevansi GWO untuk konteks IoT dan \textit{fog computing} diperkuat oleh penelitian terbaru pada \textit{task schedulling} dan \textit{resource optimization}. \textcite{alsadie_modified_2025} mengembangkan TS-GWO untuk \textit{task schedulling} IoT pada lingkunga \textit{fog-cloud} dan menunjukkan bahwa modifikasi GWO dapat meningkatkan kemampuan eksplorasi dan eksploitasi dalam menemukan solusi penjadwalan yang efisien. \textcite{alsamarai_improved_2024} juga mengusulkan PC-GWO untuk optimasi performa dan biaya pada lingkungan \textit{fog-cloud} dengan tujuan menurunkan \textit{makespan}, konsumsi energi, dan biaya komputasi. Meskipun penelitian-penelitian tersebut menunjukkan efektivitas GWO pada optimasi \textit{fog/edge}, penerapannya masih dominan pada \textit{task schedulling} dan \textit{resource allocation} dan bukan pada konfigurasi keamanan ABAC seperti TTL atribut, ukuran \textit{cache}, dan pemilihan atribut yang boleh di-\textit{cache}.

berdasarkan hal tersebut, metaheuristik digunakan karena konfigurasi \textit{cache} dan TTL pada ABAC \textit{edge} sulit ditentukan secara manual, terutama ketika tujuan optimasi mencakup performa dan keamanan secara bersamaan. Konfigurasi yang optimal tidak hanya harus menurunkan \textit{average latency, p95 latency,} dan \textit{PIP calls}, tetapi juga harus menekan \textit{stale cache hit, false allow, false deny,} dan \textit{revocation false allow}. Oleh karena itu, penggunaan GWO dalam penelitian ini diarahkan untuk mengisi celah antara optimasi performa \textit{fog/edge} yang telah banyak diteliti dan optimasi parameter keamanan ABAC yang masih relatif terbuka.